\section{Local updates}
\label{sec:updates}

The local updates are in \file{Prog/gpu\_update\_mod.F90} (633 lines: the loop over vertices, the ratio, the rank-one path) and \file{Prog/gpu\_lowrank\_mod.F90} (1723 lines: delayed and submatrix schemes, block mode, the asynchronous Intel path). This was the hardest part of the port and the part where the algorithm, not the kernel, decided the speed.

\subsection{Why the rank-one update is slow on a GPU}

Each proposal needs $G_{pp}$ of the current $G$, which depends on all earlier acceptances. A device implementation of Eq.~\eqref{eq:rank1} must, per proposal, (i) copy $G_{pp}$ to the host (or decide on the device), (ii) let the host decide, (iii) launch the update kernel. Each of these is a synchronous operation with a fixed cost (\emph{latency}) of several microseconds, independent of $N$. Measured costs:

\begin{table}[h]
\centering\footnotesize
\caption{Measured cost of single device operations.}
\label{tab:latency}
\begin{tabular}{@{}llr@{}}
\toprule
Platform, operation & Source & $\mu$s \\
\midrule
A30, empty OpenMP target region, synchronous & \file{BENCHMARKS.md}, ticket 102 & 12.4 \\ % BENCHMARKS.md
A30, empty OpenMP target, \texttt{nowait} & same & 11.0 \\
A30, empty OpenACC kernel, \texttt{async} & same & 2.7 \\
A30, one proposal with host decision & same & 15.2 \\
A30, D$\to$H$\to$D round trip, $N\le400$ & \file{ADR-001} & 28--34 \\ % ADR-001 context
A30, rank-1 update $N=144/400/1024$ (device) & \file{BENCHMARKS.md} & 4.7 / 6.6 / 31.6 \\
A30, rank-$k$ ($k=128$) per update, same $N$ & same & 0.21 / 0.37 / 1.25 \\
Max 1550, empty synchronous target & \file{SYNC.md} & 6.1 \\ % SYNC.md
Max 1550, small kernel & same & 9--12 \\
Max 1550, ZGEMM $1024\times32\times32$ / $1024\times1024\times32$ & same & 17.7 / 55 \\
Max 1550, D$\to$H 16~B / 512~KB & same & 2.8 / 350 \\
Max 1550, SYCL submit, asynchronous & \file{p3/STATUS.md} & 4.3 \\ % p3 STATUS phase 3b
MI210, GCC (libgomp) target region with parallel loop & \file{alf\_hip\_kernels.hip.cpp:2} & $\approx$150 \\
MI210, HIP launch plus synchronisation & same & $\approx$15 \\
\bottomrule
\end{tabular}
\end{table}

With $2N\Lt$ proposals per sweep and about 30~$\mu$s per proposal, a sweep at $L=24$ ($N=576$, $\Lt=240$) would spend $2\cdot576\cdot240\cdot30~\mu{\rm s}\approx8$~s in latency alone, before any arithmetic. The rank-$k$ row of Table~\ref{tab:latency} shows the way out: applying $k=128$ updates as one matrix product costs 25 to 40 times less per update. Two schemes that collect updates, the delayed and the submatrix update, exist as CPU code in \file{Lowrank\_update\_mod.F90} (taken into the branch on 2026-11-06 from a fork branch; the upstream master \texttt{521ddfca} has no such module). The port implements both on the device; the submatrix scheme in \emph{block mode} is the production path.

\subsection{Common part: the ratio and the decision}

For all schemes the host does the same arithmetic as \texttt{Upgrade2} (\file{gpu\_update\_mod.F90:361-465}). Only single-site diagonal vertices are supported ($N_{\rm non\text{-}zero}\le1$). With $z_1=g\,(\eta_{\rm new}-\eta_{\rm old})$, the vertex value $E_1$ and $\alpha$:
\begin{align}
\Delta &= e^{z_1E_1}-1, \qquad d = 1+\Delta\,(1-G_{pp}),\\
R_f &= d\;e^{z_1\alpha},\qquad R = \Bigl(\prod_f R_f\Bigr)^{N_{\rm SUN}}\frac{\gamma_{\rm new}}{\gamma_{\rm old}},\\
W &= S_0\,T_0\,\left|\Re(\Phi R)/\Re\Phi\right|,\qquad\text{accept if } W>\texttt{ranf\_wrap()}.
\end{align}
The random numbers are drawn on the host in the same order as on the CPU (the $T_0$ test draw, then the Metropolis draw). The only input that depends on the scheme is $G_{pp}$. On acceptance the update of $G$ uses $z=\Delta/d$.

\subsection{Standard (rank-one) update}

$G_{pp}$ is copied D$\to$H (16 bytes). On acceptance two kernels run per flavour: \texttt{k\_upd\_prep} sets $x=G_{:,p}$ and $y=-G_{p,:}$ with $y_p=1-G_{pp}$, and \texttt{k\_upd\_geru} performs $G_{ij}\leftarrow G_{ij}-z\,x_iy_j$, which is Eq.~\eqref{eq:rank1} (\file{gpu\_update\_mod.F90:301-334}). The new field is written into \texttt{nsigma\_d} by a 16-byte H$\to$D. The complex scalar $z$ is passed as two reals: nvfortran 26.9 dropped the imaginary part of a complex scalar argument inside a target region, which gave errors of $10^{-3}$ to 1 in $G$ (\file{DISCREPANCIES.md}). % DISCREPANCIES.md

\subsection{Delayed update}

\paragraph{Representation.} Per flavour $G = G_0 - XY^T$ with $G_0=$ \texttt{G\_d} and $X,Y\in\mathbb{C}^{N\times\Nd}$ on the device; $j$ columns are in use (\file{gpu\_lowrank\_mod.F90:39-45}).
\paragraph{Proposal.} $G_{pp}=(G_0)_{pp}-\sum_{l\le j}X_{pl}Y_{pl}$, computed by a single-thread kernel; two reals go D$\to$H.
\paragraph{Accept.} One kernel computes the current column and row of $G$ and appends them:
\begin{align}
c &= G_0(:,p) - X\,Y(p,:)^T = G(:,p),\qquad r = G_0(p,:) - X(p,:)\,Y^T = G(p,:),\\
X(:,j{+}1) &= z\,c,\qquad Y(:,j{+}1) = e_p - r^T .
\end{align}
Then $G_0-XY^T$ after the append equals $G - z\,G(:,p)\,(e_p^T-G(p,:))$, which is Eq.~\eqref{eq:rank1}. By induction, after $j$ acceptances $G_0-XY^T$ equals $j$ sequential rank-one updates; only the summation order differs.
\paragraph{Wrap.} The vertex wrap $G\to LGR$ with diagonal $L$, $R$ that differ from the identity only at site $r$ is applied to $G_0$ (row $r$ times $z_a$, column $r$ times $z_b$) and to the factors, $X(r,:)\leftarrow z_aX(r,:)$, $Y(r,:)\leftarrow z_bY(r,:)$. This is exact: $LGR=LG_0R-(LX)(RY)^T$.
\paragraph{Flush.} $G_0\leftarrow G_0-X(:,1{:}j)\,Y(:,1{:}j)^T$, one \texttt{la\_gemm} of size $N\times N\times j$, when $j=\Nd$, before a vertex with $N>1$ and at the end of every \texttt{WRAPGRUP}/\texttt{WRAPGRDO}. The hopping wraps, stabilisation and measurements always see a flushed $G$.

The delayed scheme removes the $N^2$ work per acceptance but still needs one round trip per proposal. A ``fast'' variant handles both flavours in one launch per site and defers the field write and the next wrap into the arguments of the next launch (\file{gpu\_lowrank\_mod.F90:699-733, 909-1055}).

\subsection{Submatrix update}

\paragraph{Representation.} Per flavour, with the set $S=\{s_1,\ldots,s_k\}$ of distinct pending sites ($k\le\Nd$) and a small matrix $M\in\mathbb{C}^{k\times k}$ kept on the \emph{host} (\file{gpu\_lowrank\_mod.F90:46-51}):
\begin{equation}
\label{eq:subrep}
G = G_0 - G_0(:,S)\;M\;\bigl(E_S^T-G_0(S,:)\bigr),
\end{equation}
where $E_S$ is the $N\times k$ matrix of unit columns $e_{s_m}$. Between flushes $G_0$ is not changed by updates.

\paragraph{Proposal.} Let $J=e_{{\rm slot}(p)}\in\mathbb{C}^k$ if $p\in S$ and $J=0$ otherwise. From Eq.~\eqref{eq:subrep}, since $(E_S^T-G_0(S,:))(:,p)=J-G_0(S,p)$,
\begin{equation}
G(:,p) = G_0(:,p) - G_0(:,S)\,w,\qquad w = M\,\bigl(J-G_0(S,p)\bigr),\qquad G_{pp} = (G_0)_{pp}-G_0(p,S)\,w .
\end{equation}
This needs only $(G_0)_{pp}$, $G_0(S,p)$ and $G_0(p,S)$: $2k+1$ numbers. The code (\texttt{gpu\_lr\_gpp}) computes $w$ and $G_{pp}$ on the host in $O(k^2)$ and caches $w$ and $G_0(p,S)$ for the accept.

\paragraph{Accept (host only).} If $p\notin S$, $S$ is padded with $p$ and $M$ with a zero row and column ($k'=k+1$, slot $i=k'$); otherwise $i={\rm slot}(p)$ and $k'=k$. Define
\begin{equation}
a = e_i - \begin{pmatrix} w\\0\end{pmatrix},\qquad b^T = e_i^T + \bigl(G_0(p,S)\,M,\;0\bigr),\qquad M\leftarrow M + z\,a\,b^T .
\end{equation}
No device work at all (\file{gpu\_lowrank\_mod.F90:649-696}).

\paragraph{Equivalence with the rank-one update.} With the padded set $S'$ and $C=G_0(:,S')$, $Q=E_{S'}^T-G_0(S',:)$: since $p\in S'$, $G_0(:,p)=Ce_i$, hence $G(:,p)=C(e_i-[w;0])=Ca$. Further $e_p^T-G_0(p,:)$ is row $i$ of $Q$, and $G_0(p,S)MQ_S = [G_0(p,S)M,0]\,Q$, hence $e_p^T-G(p,:)=b^TQ$. The Sherman--Morrison step becomes
\begin{equation}
G' = G - z\,G(:,p)\bigl(e_p^T-G(p,:)\bigr) = G_0 - C\,(M+z\,ab^T)\,Q ,
\end{equation}
which is the update of $M$ above. The formula covers new, repeated and overlapping sites in one expression, as on the CPU (\file{Lowrank\_update\_mod.F90:1187-1198}).

\paragraph{Wrap.} $G\to LGR$ with $L$, $R$ diagonal and different from $I$ only at site $r$: $G_0\leftarrow LG_0R$ on the device, and if $r\in S$, row ${\rm slot}(r)$ of $M$ times $z_a$ and column ${\rm slot}(r)$ times $z_b$ on the host. Writing $G_0'=LG_0R$ one finds $LG_0(:,S)=G_0'(:,S)R_S^{-1}$ and $(E_S^T-G_0(S,:))R=R_SE_S^T-L_S^{-1}G_0'(S,:)$. The representation is preserved with $M'=L_SMR_S$ exactly when $R_S=L_S^{-1}$, i.e.\ $z_az_b=1$. For the vertices supported this holds: $z_a=E(s)$ and $z_b=E(-s)$ with $\eta(-s)=-\eta(s)$. The code does not assert it.

\paragraph{Flush.} With $C=G_0(:,S)$ and $Q^T=E_S-G_0(S,:)^T$ gathered by one kernel into \texttt{X\_d} and scratch,
\begin{equation}
Y = Q^T M^T \;(N\times k\times k),\qquad G_0\leftarrow G_0 - C\,Y^T\;(N\times N\times k),
\end{equation}
two \texttt{la\_gemm} calls (\file{gpu\_lowrank\_mod.F90:480-497}). $C\,Y^T=CMQ$ is the correction of Eq.~\eqref{eq:subrep}.

\paragraph{$\Nd$.} If \texttt{N\_delay}$\ge1$ it is used as given; otherwise the static rule of the CPU (\texttt{Lowrank\_static\_N\_delay}) applies, and on the GPU the submatrix scheme additionally uses at least $\min(32,n_{\rm nz})$: ``larger blocks amortise the device round trips'' (\file{gpu\_lowrank\_mod.F90:253}). \texttt{N\_delay=-1} (automatic tuning on the CPU) is treated like 0 on the GPU.

\subsection{Block mode}

The submatrix scheme as described still needs one gather and one D$\to$H per proposal. Block mode removes them (\file{gpu\_lowrank\_mod.F90:1253-1260, 1431-1720}). It is active when the scheme is submatrix and the ``fast'' conditions hold: $N_{\rm FL,eff}\le2$, every vertex a single-site operator, and the sites of different vertices distinct (\texttt{fast\_check}).

\paragraph{Idea.} Between flushes $G_0$ changes only by the diagonal wrap scalings, and a proposal at $p$ needs only $(G_0)_{pp}$, $G_0(p,S)$, $G_0(S,p)$ with $S\subset B$, the sites of the block. So the host can copy the principal block $G_0(B,B)$ once, apply each wrap to its copy, and decide all proposals of the block without the device.

\paragraph{A block} is $b=\min(\Nd,\text{remaining vertices of the slice})$ consecutive vertices in sweep order (ascending in \texttt{WRAPGRUP}, descending in \texttt{WRAPGRDO}); blocks do not cross slices. Since the sites are distinct, $k\le b\le\Nd$, so no capacity flush can happen inside a block.

\paragraph{Sequence of one block.}
\begin{enumerate}
\item \texttt{blk\_begin}: the host builds the list of block sites; H$\to$D of the list; kernel \texttt{k\_blk\_gather} copies $G_0(B,B)$ of both flavours into \texttt{gbd}; D$\to$H of \texttt{gbd} ($32\Nd^2$ bytes).
\item For each vertex, on the host only: wrap of the host block (row $a$ times $z_a$, column $a$ times $z_b$; the factors are recorded) and of $M$; proposal from the host block by the formulas above; accept by the host update of $M$. On the way up the wrap precedes the proposal; on the way down it follows the update with the final field.
\item \texttt{blk\_end}: the host packs the recorded factors and $M$ into \texttt{cbuf}; H$\to$D of \texttt{cbuf} and of $S$; kernel \texttt{k\_blk\_scale} applies all row factors and then all column factors of the block sites to $G_0$ (both flavours, in the order of the two-kernel version); per flavour the flush (gather plus two GEMMs); H$\to$D of the new fields \texttt{nsigma\_d(vmin:vmax,t)}.
\end{enumerate}
A flush with $k=1$ is padded to $k=2$ with a duplicate site and a zero row and column of $M$ (\file{gpu\_lowrank\_mod.F90:1528-1532}); the extra site contributes nothing to $CMQ$ because its row and column of $M$ are zero, so the result is the same as a rank-one flush up to rounding. Neither the code nor the commit that introduced it (\texttt{e70a3061}) gives a reason \todo{reason for the k=1 to k=2 padding in blk\_end; possibly to avoid a library special case for $k=1$, not documented}.

At $L=24$, $\Nd=32$ one sweep has 8\,640 blocks (18 per slice $\times$ 480 slices) with 8\,640 gathers, 17\,280 flush gathers, 8\,640 scale kernels, 34\,560 GEMMs and 43\,200 copies (\file{p7/NOTES.md}, line 94). % p7/NOTES.md
The flush rank is small: $k$ of 10 to 14 on average (\file{p7/NOTES.md}, line 36). % p7 NOTES

\paragraph{Effect.} On the A30 at $N=400$ the per-proposal path cost 117~$\mu$s per proposal (5.0 of 7.9~s); block mode brought this part to 0.53~s, ten times less (\file{DELAYED\_SUBMATRIX.md}, line 52). % DELAYED_SUBMATRIX.md
The submatrix scheme with block mode is 2.4 to 5 times faster than the standard update and 1.4 to 2.1 times faster than the delayed scheme at the same $\Nd$ ($L=16$ to 32, A30, \file{TUNING\_DEFECT.md}). % TUNING_DEFECT.md

On nvfortran 26.9 the routine \texttt{blk\_end} is compiled with \texttt{!pgi\$r opt=1}: at \texttt{-O2}/\texttt{-O3} the compiler produced a wrong $G$ after the first stabilisation (\file{gpu\_lowrank\_mod.F90:1496-1499}).

\begin{figure}[h]
\centering
\begin{tikzpicture}[x=0.95cm,y=0.8cm]
\node[anchor=east,font=\small] at (0,2) {host};
\node[anchor=east,font=\small] at (0,0) {device};
\draw[gray!50] (0,2.5)--(15.5,2.5); \draw[gray!50] (0,1)--(15.5,1); \draw[gray!50] (0,-0.5)--(15.5,-0.5);
% block begin
\filldraw[fill=gray!20] (0.2,1.6) rectangle (1.0,2.4); \node[font=\scriptsize] at (0.6,2.0) {sites};
\draw[arr] (1.0,1.7) -- (1.3,0.4);
\filldraw[fill=cA30!30] (1.3,-0.4) rectangle (2.3,0.4); \node[font=\scriptsize] at (1.8,0) {gather};
\draw[arr] (2.3,0.35) -- (2.7,1.65); \node[font=\scriptsize,red!70!black,anchor=east] at (2.45,1.0) {wait};
% host proposals
\foreach \i in {0,...,7} { \filldraw[fill=gray!35] (2.8+\i*0.8,1.6) rectangle (3.5+\i*0.8,2.4); }
\node[font=\scriptsize,align=center] at (6.0,2.8) {$b$ vertices: wrap block copy, propose, decide, update $M$};
\draw[arr] (9.2,1.7) -- (9.5,0.4);
\filldraw[fill=cA30!30] (9.5,-0.4) rectangle (10.5,0.4); \node[font=\scriptsize] at (10.0,0) {scale};
\filldraw[fill=cA30!50] (10.5,-0.4) rectangle (11.2,0.4); \node[font=\scriptsize] at (10.85,0) {fg};
\filldraw[fill=cA30!70] (11.2,-0.4) rectangle (12.6,0.4); \node[font=\scriptsize] at (11.9,0) {2 GEMM};
\filldraw[fill=cA30!50] (12.6,-0.4) rectangle (13.3,0.4); \node[font=\scriptsize] at (12.95,0) {fg};
\filldraw[fill=cA30!70] (13.3,-0.4) rectangle (14.7,0.4); \node[font=\scriptsize] at (14.0,0) {2 GEMM};
\node[font=\scriptsize,align=left,anchor=west] at (9.5,-0.95) {flavour 1 \hspace{1.6cm} flavour 2};
\filldraw[fill=gray!20] (14.6,1.6) rectangle (15.6,2.4); \node[font=\scriptsize] at (15.1,2.0) {next};
\end{tikzpicture}
\caption{One block of the submatrix update in block mode (schematic, not to scale). Host work is grey, device work green; ``fg'' is the flush gather. The only unavoidable wait is the copy of $G_0(B,B)$ to the host at the beginning of the block, because the next gather depends on the previous flush. In the synchronous path every arrow is also a wait.}
\label{fig:block}
\end{figure}

\subsection{Synchronisation points per block}

In the synchronous OpenMP path every \texttt{target} region and every \texttt{target update} returns only when it is complete. In addition \texttt{alf\_la} calls run asynchronously on NVIDIA (cuBLAS returns before the GEMM is done), so two explicit \texttt{la\_sync} calls are needed. The sequence per block (\file{gpu\_lowrank\_mod.F90:1443-1592}):
\begin{enumerate}
\item \texttt{la\_sync} in \texttt{blk\_begin}: the previous flush (asynchronous library calls) must be complete before new uploads overwrite buffers it reads.
\item \texttt{target update to(bsd)}: the block sites.
\item \texttt{k\_blk\_gather}.
\item \texttt{target update from(gbd)}: the host needs $G_0(B,B)$ before it can decide anything. This is the one wait that cannot be removed.
\item \texttt{la\_sync} in \texttt{blk\_end}: ``the previous block's flush (cuBLAS) is asynchronous: it must be done before cbuf / bS\_d are overwritten''.
\item \texttt{target update to(cbuf, bS\_d)}.
\item \texttt{k\_blk\_scale}.
\item Per flavour: \texttt{k\_sub\_flush\_gather} and two \texttt{la\_gemm}. On Intel the dispatched GEMM includes an implicit wait (\file{SYNC.md}).
\item \texttt{target update to(nsigma\_d(vmin:vmax,t))}.
\end{enumerate}
On the A30 a profile of a production layout showed 72\,\% of host time in \texttt{cuStreamSynchronize}, with about 4\,300 blocks and 44\,000 synchronisations per sweep (\file{PP\_GPU.md}, $L=24$, 9 ranks). % PP_GPU.md
In the single-chain profile at $L=32$, 44\,\% of the kernel time was in 32\,338 small ZGEMM flushes of 52~$\mu$s each (\file{TUNING\_DEFECT.md}). % TUNING_DEFECT.md

\subsection{Launch-latency budget of a block}

The Intel team measured where the 167~$\mu$s of wall time of one block go at $L=24$ with V7 (\file{p7/NOTES.md}, line 97): % p7 NOTES
submission of the asynchronous operations 91~$\mu$s, the D$\to$H copy 22~$\mu$s (30 at $L=32$), the hopping 32~$\mu$s, host arithmetic about 20~$\mu$s; the device executes the block in about 65~$\mu$s. The update phase is therefore ``host-submission bound'': the cost of telling the device what to do exceeds the work it does. The per-operation submission costs were 4.4~$\mu$s (gather), 4.5 (flush gather), 4.4 (scale), 9.7 (GEMM) and 9.0 (copy), 0.88~s per sweep in total, 61\,\% of the update phase (\file{p7/NOTES.md}, lines 88-97).

Before the staging buffer setting the D$\to$H of the 32~KB block cost 380~$\mu$s on Intel; with \env{LIBOMPTARGET\_LEVEL\_ZERO\_STAGING\_BUFFER\_SIZE=1048576} it costs 22~$\mu$s, and the update phase went from 4.3 to 1.9~s per sweep (\file{p3/STATUS.md}; \file{ISSUE\_INTEL\_VS\_A30.md}, line 43). % p3 STATUS, ISSUE line 43
``One environment variable gained more than two code phases'' (\file{GPU\_PLAN\_V2.md}, section 2).

\subsection{The asynchronous block path on Intel}
\label{sec:asyncblk}

On NVIDIA, cuBLAS calls are asynchronous by default. On Intel, ifx executes OpenMP \texttt{nowait} target tasks undeferred (97 to 103\,\% of the synchronous cost, \file{p3/STATUS.md}), so OpenMP cannot overlap host and device. The asynchronous block path (\texttt{ALFSH\_ASYNCBLK=1}, \file{gpu\_lowrank\_mod.F90:1328-1347, 1605-1617}; \file{alf\_la\_sycl\_blk.cpp}) bypasses OpenMP for the block operations:
\begin{itemize}
\item It obtains the SYCL queue that the OpenMP runtime uses (\texttt{alfsh\_get\_queue}, through the interop object of \texttt{alf\_la}) and refuses it if it is not \emph{in-order}. ``The queue must be in-order: that, and nothing else, orders the calls'' (\file{alf\_la\_sycl\_blk.cpp:4}).
\item The gather, scale and flush-gather kernels are SYCL ports of the OpenMP kernels with the same expressions and order; the GEMM is oneMKL's \texttt{column\_major::gemm} with the same shapes and scalars; copies are \texttt{queue.memcpy}. All are submitted without waiting. The gathered block lands in host USM memory (memory that the device can write directly).
\item The only wait per block is \texttt{ab\_drain} in \texttt{blk\_begin} after the gather: ``the one unavoidable wait of a block: the host needs G0(B,B)''. \texttt{blk\_end} leaves the block pending; the next \texttt{blk\_begin}, the flush at the end of the slice, or any other device access drains it.
\item The SYCL code is built as a shared library (\file{libalfsb.so}) because SYCL kernels in an object linked by ifx are not found at run time (``No kernel named \ldots'').
\end{itemize}
Results are bitwise identical to the synchronous path. Planted tests: a missing wait gives an ``Unstable matrix'' abort; a missing drain at the end of the slice changes nothing, because the in-order queue orders the OpenMP kernels too (\file{p3/STATUS.md}). Measured effect: $L=24$ 7.05--7.08$\to$6.53--6.59~s per sweep, $L=32$ 23.2$\to$22.0--22.4~s, per card 0.568$\to$0.655--0.688 sweeps/s (\file{ISSUE\_INTEL\_VS\_A30.md}, section 2). % ISSUE section 2

\dis{The release status of \texttt{ALFSH\_ASYNCBLK} differs between notes: \file{RUNTIME\_SWITCHES.md} (line 123) lists it as production, \file{SPEEDUPS.md} (line 21) and \file{NOT\_IMPLEMENTED.md} (line 118) as not yet released. The measured gain is +3\,\% per node in \file{p3/STATUS.md} and +15 to 20\,\% per card in \file{ISSUE\_INTEL\_VS\_A30.md}; the jobs differ.}

\subsection{What remains: the flush}

In production layouts the flush GEMMs are 47\,\% of the kernel time (A30, 9 ranks, \file{PP\_GPU.md}). % PP_GPU.md
They are rank-$k$ updates with $k\approx10$ to 20, bound by memory bandwidth ($N^2$ elements read and written per flush for $2Nk$ new information). The ranked candidate is a two-level submatrix update: one flush per slice of rank about 190 instead of one per block (\file{SPEEDUPS.md}, section 4). It is estimated at weeks of work and not started.

\dis{The header of \file{gpu\_lowrank\_mod.F90} (lines 34, 52-53) calls the module a ``test branch'' where ``numerical agreement is not the goal''. It is the production path, and \file{DELAYED\_SUBMATRIX.md} reports bin-by-bin agreement. The header also describes only the per-proposal gather; block mode is documented only at line 1253.}
